LiionTR is organized around a modular formulation in which physical
models are separated from numerical integration. The objective is to
represent lithium-ion battery thermal runaway using a set of independent
but coupled subsystems describing reaction kinetics, heat generation,
gas generation, pressure evolution, and venting.

The state of a zero-dimensional LiionTR simulation can be written
conceptually as

\begin{equation}
    \mathbf{y}
    =
    \left[
        T,
        \alpha_1,
        \ldots,
        \alpha_N,
        n_1,
        \ldots,
        n_M
    \right]^{\mathrm{T}},
\end{equation}

where $T$ is the lumped cell temperature, $\alpha_i$ is the conversion
of thermal runaway reaction $i$, and $n_s$ is the amount of gas species
$s$ contained in the cell free volume.

The governing system may therefore be expressed as

\begin{equation}
    \frac{d\mathbf{y}}{dt}
    =
    \mathbf{F}
    \left(
        t,
        \mathbf{y}
    \right),
\end{equation}

where the right-hand side is assembled from the thermal, reaction,
gas-generation, pressure, and vent-flow models.


\subsection{Thermal state}

The cell temperature evolves according to an energy balance,

\begin{equation}
    C_{\mathrm{th}}
    \frac{dT}{dt}
    =
    \dot{Q}_{\mathrm{gen}}
    -
    \dot{Q}_{\mathrm{loss}},
    \label{eq:general_energy_balance}
\end{equation}

where $C_{\mathrm{th}}$ is the effective cell thermal capacity,
$\dot{Q}_{\mathrm{gen}}$ is the total internal heat-generation rate, and
$\dot{Q}_{\mathrm{loss}}$ is the heat-transfer rate from the cell to its
environment.


\subsection{Reaction states}

Each thermal runaway reaction is represented using a scalar conversion,

\begin{equation}
    0
    \leq
    \alpha_i
    \leq
    1.
\end{equation}

The corresponding reaction rate is represented by

\begin{equation}
    \frac{d\alpha_i}{dt}
    =
    k_i(T)
    f_i
    \left(
        \alpha_i,
        \mathbf{\alpha}
    \right),
    \label{eq:reaction_progress_general}
\end{equation}

where $k_i(T)$ describes the intrinsic temperature dependence and
$f_i$ describes the dependence on reaction progress and, when required,
the state of other reactions.

The complete reaction-state vector is

\begin{equation}
    \boldsymbol{\alpha}
    =
    \left[
        \alpha_1,
        \alpha_2,
        \ldots,
        \alpha_N
    \right]^{\mathrm{T}}.
\end{equation}


\subsection{Heat generation}

For a reaction occupying an effective cell mass fraction $w_i$ and
releasing a specific enthalpy $\Delta H_i$, the mass-specific
heat-generation rate is written as

\begin{equation}
    \dot{q}_i
    =
    w_i
    \Delta H_i
    \frac{d\alpha_i}{dt}.
\end{equation}

The total heat-generation rate for a cell of mass $m_{\mathrm{cell}}$
is then

\begin{equation}
    \dot{Q}_{\mathrm{gen}}
    =
    m_{\mathrm{cell}}
    \sum_{i=1}^{N}
    \dot{q}_i.
    \label{eq:total_heat_generation}
\end{equation}


\subsection{Gas-generation states}

When gas generation is enabled, each gas species is represented by a
mole inventory $n_s$.

Before vent opening,

\begin{equation}
    \frac{dn_s}{dt}
    =
    \dot{n}_{s,\mathrm{gen}},
\end{equation}

where $\dot{n}_{s,\mathrm{gen}}$ is the molar gas-generation rate.

After vent opening,

\begin{equation}
    \frac{dn_s}{dt}
    =
    \dot{n}_{s,\mathrm{gen}}
    -
    \dot{n}_{s,\mathrm{vent}}.
    \label{eq:vented_species_balance}
\end{equation}


\subsection{Pressure state}

The current reduced-order formulation evaluates internal gas pressure
using the ideal-gas equation,

\begin{equation}
    P
    =
    \frac{
        n_{\mathrm{tot}} R T
    }{
        V_{\mathrm{free}}
    },
    \label{eq:ideal_pressure_framework}
\end{equation}

where

\begin{equation}
    n_{\mathrm{tot}}
    =
    \sum_{s=1}^{M}
    n_s.
\end{equation}

The current formulation assumes that the gas temperature is equal to the
lumped cell temperature.


\subsection{Model coupling}

The principal physical couplings represented in LiionTR are

\begin{enumerate}
    \item temperature controls thermal runaway reaction rates;
    \item thermal runaway reactions release heat;
    \item heat generation modifies the cell temperature;
    \item reaction progress can generate gas;
    \item the gas inventory and temperature determine internal pressure;
    \item internal pressure can trigger vent opening;
    \item vent flow removes gas from the internal inventory.
\end{enumerate}

The reaction-thermal coupling forms the fundamental positive feedback
mechanism responsible for thermal runaway.

A conceptual representation of the current LiionTR architecture is shown
in Figure~\ref{fig:framework_architecture}.

\begin{figure}[htbp]
    \centering
    \begin{tikzpicture}[
        node distance=1.8cm and 2.4cm,
        block/.style={
            rectangle,
            rounded corners,
            draw,
            align=center,
            minimum width=3.0cm,
            minimum height=1.0cm
        },
        line/.style={
            -{Latex[length=3mm]},
            thick
        }
    ]

    \node[block] (kinetics) {Reaction\\kinetics};
    \node[block, below=of kinetics] (progress) {Reaction\\progress};
    \node[block, below=of progress] (heat) {Heat\\generation};
    \node[block, below=of heat] (thermal) {Thermal\\model};

    \node[block, right=of progress] (gasgen) {Gas\\generation};
    \node[block, below=of gasgen] (inventory) {Gas\\inventory};
    \node[block, below=of inventory] (pressure) {Pressure};
    \node[block, below=of pressure] (venting) {Venting};

    \draw[line] (kinetics) -- (progress);
    \draw[line] (progress) -- (heat);
    \draw[line] (heat) -- (thermal);

    \draw[line] (progress) -- (gasgen);
    \draw[line] (gasgen) -- (inventory);
    \draw[line] (inventory) -- (pressure);
    \draw[line] (pressure) -- (venting);

    \draw[line] (thermal.east) -- ++(1.2,0) |- (pressure.west);
    \draw[line] (pressure.west) -- ++(-1.2,0) |- (thermal.east);

    \end{tikzpicture}
    \caption{
        Conceptual architecture of the current LiionTR framework.
        Reaction kinetics determine reaction progress, which drives both
        heat generation and gas generation. Heat generation affects the
        thermal model, while gas generation affects the gas inventory and
        therefore the internal pressure. Pressure may eventually trigger
        venting.
    }
    \label{fig:framework_architecture}
\end{figure}


\subsection{Separation of physical models and state integration}

A central design principle of LiionTR is that physical model objects do
not advance their own numerical state.

Instead, the numerical solver owns the complete state vector and passes
the current state to the physical models whenever the right-hand side is
evaluated.

This design provides several advantages:

\begin{itemize}
    \item deterministic evaluation of the governing equations;
    \item compatibility with standard ordinary differential equation
          solvers;
    \item reduced hidden coupling between physical components;
    \item independent verification of individual physical models;
    \item easier implementation of event handling and future alternative
          numerical solvers.
\end{itemize}


\subsection{Scope of the present framework}

The current implementation focuses on reduced-order, zero-dimensional
thermal runaway modeling.

The framework is intended to evolve toward more comprehensive physical
descriptions, including thermochemical equilibrium, detailed
electrochemical state information, spatially resolved heat transfer,
cell-to-cell propagation, structural failure, and external vent-jet or
combustion models.

Cantera is considered as a future thermochemical backend, while PyBaMM
may provide electrochemical states and normal-operation battery physics.
LiionTR is intended to remain the orchestration layer responsible for
thermal runaway and battery-safety-specific coupling.